\documentclass[11pt]{article}
\usepackage{amsmath,amssymb,amsthm,hyperref,xurl}
\setlength{\emergencystretch}{3em}
\renewcommand{\labelitemi}{$\bullet$}
\usepackage[margin=1in]{geometry}
\newcommand{\N}{\mathbb N}
\newcommand{\Imp}{\mathsf I}
\newcommand{\Rel}{\mathsf R}
\newcommand{\All}{\mathsf A}
\newcommand{\Eq}{\mathsf E}
\newcommand{\Neg}{\mathsf N}
\newcommand{\Ex}{\mathsf X}
\newcommand{\J}{\mathsf J}
\newtheorem{theorem}{Theorem}
\newtheorem{lemma}{Lemma}
\title{A one-stack polynomial core for ZFC refutations}
\author{Technical construction and pinned-source extraction}
\date{22 October 2026}
\begin{document}
\maketitle

\paragraph{Scope.}
This note gives an explicit polynomial \emph{local} transition relation whose
finite accepting paths are equivalent to finite ZFC refutations. It removes
the tape and enumeration layers of a certified source calculus. Compressing
an unbounded finite path to a fixed number of Diophantine variables is a
separate step; this note does not claim that a variable-length family of
equations is already the final fixed-arity Diophantine system.
All variables below range over \(\N=\{0,1,2,\ldots\}\).

\section{The extracted numerical calculus}

Define the following integer polynomials, with abbreviations expanded by
composition, not treated as additional function symbols:
\begin{align*}
 \Imp(x,y)&=2(x+y)(x+y+1)+4y,\\
 \Rel(x,y)&=\Imp(x,y)+1,&
 \All(x,y)&=\Imp(x,y)+2,&
 \Eq(x,y)&=\Imp(x,y)+3,\\
 \Neg(x)&=\Imp(x,1),&
 \Ex(x,y)&=\Neg(\All(x,\Neg(y))).
\end{align*}
These are exactly the source codes \(4P(x,y)+j\), where
\(P(x,y)=(x+y)(x+y+1)/2+y\). In particular, there is no division in
the displayed integer-polynomial constructors.

The eight logical rows of the calculus are
\begin{align*}
 r_1(a,b)&=\Ex(a,\Eq(a,b)),\\
 r_2(a,b,c)&=\Imp(\All(a,\Imp(b,c)),
                         \Imp(\All(a,b),\All(a,c))),\\
 r_3(a,b,c)&=\Imp(\Eq(a,b),\Imp(\Rel(c,a),\Rel(c,b))),\\
 r_4(a,b,c)&=\Imp(\All(a,\All(b,c)),\All(b,\All(a,c))),\\
 r_5(a,b,c)&=\Imp(\Eq(c,a),\Imp(\Rel(c,b),\Rel(a,b))),\\
 r_6(a,b)&=\Imp(\Ex(a,\All(a,b)),b),\\
 g(a,b,c)&=\Imp(\Eq(c,a),
                 \Imp(\Eq(c,b),\All(\Rel(a,b),\Eq(a,b)))),\\
 r_8(a,b,c,u)&=\Imp(\Imp(\Imp(c,u),a),
                        \Imp(\Imp(a,c),\Imp(b,c))).
\end{align*}
For an unambiguous short definition of the seed bundle put
\begin{align*}
 Q(b,k,m)&=\All\bigl(b,\Imp(m,
       \Neg(\All(b,\Imp(m,\Neg(\Rel(b,k))))))\bigr),\\
 V(a,b,k)&=\All(k,\Imp(\Rel(k,a),\Rel(k,b))),\\
 C(a,b,k,m)&=\All\bigl(b,\Imp(\Rel(b,k),
    \Neg(\All(a,\Imp(m,
      \Neg(\All(a,\Imp(V(a,b,k),Q(b,k,m))))))))\bigr),\\
 K(a,b,c)&=\All(c,\Imp(C(a,b,c,c),\Neg(Q(b,c,c)))),\\
 R_V(a,b,c)&=\Imp(\Rel(a,b),\All(c,\Rel(a,b))),\\
 B(a,b,c)&=\Neg(\Imp(R_V(a,b,c),K(a,b,c))).
\end{align*}

Let \(D\subseteq\N\) be the least set satisfying these twelve rules:
\begin{enumerate}
\item \(0\in D\).
\item If \(a\in D\) and \(\Imp(a,b)\in D\), then \(b\in D\).
\item If \(b\in D\), then \(\All(a,b)\in D\), for every \(a\).
\item Each value of \(r_1,r_2,r_3,r_4,r_5,r_6,g,r_8\) is in \(D\)
      (eight separate rules).
\item \(B(a,b,c)\in D\) whenever \(c\ne a\) and \(c\ne b\).
\end{enumerate}
Thus \(D(n)\) is exactly \texttt{KCode222.C222 n}, and falsity is the
number \(1\). A derivation means a finite tree built from these rules.

\section{A polynomial one-stack checker}

Encode the empty stack by 0, and a stack with head \(x\) and tail \(y\) by
\[
  \J(x,y)=(x+y)^2+y+1.
\]
No surjectivity onto positive integers is required.

\begin{lemma}
The map \(\J:\N^2\to\N\) is injective and positive.
\end{lemma}
\begin{proof}
For \(x+y=n\), its values lie in the interval
\([n^2+1,n^2+n+1]\), and within this interval the value determines \(y\).
The next interval begins at \((n+1)^2+1=n^2+2n+2\), which is strictly
greater than \(n^2+n+1\). Hence different diagonals have disjoint ranges,
and both coordinates are determined. Positivity is immediate.
\end{proof}

For natural \(a,b,c,u\), let the nine-entry list \(\mathcal Q\) be
\[
 \mathcal Q(a,b,c,u)=
 (0,r_1(a,b),r_2(a,b,c),r_3(a,b,c),r_4(a,b,c),
 r_5(a,b,c),r_6(a,b),g(a,b,c),r_8(a,b,c,u)).
\]
Repeated values in this list are harmless: the product below is over
its nine entries, with multiplicity. Define three sum-of-squares factors:
\begin{align*}
 F_B(s,t;a,b,c,u,v)
   &=(t-\J(B(a,b,c),s))^2\\
   &\quad+((c-a)^2-u-1)^2+((c-b)^2-v-1)^2,\\
 F_G(s,t;a,b,u)
   &=(s-\J(b,u))^2+(t-\J(\All(a,b),u))^2,\\
 F_M(s,t;a,b,u)
   &=(s-\J(\Imp(a,b),\J(a,u)))^2+(t-\J(b,u))^2.
\end{align*}
The promised single local polynomial is
\[
 \boxed{\displaystyle
 H(s,t;a,b,c,u,v)=
 \left(\prod_{q\in\mathcal Q(a,b,c,u)}(t-\J(q,s))\right)
 F_B(s,t;a,b,c,u,v)F_G(s,t;a,b,u)F_M(s,t;a,b,u).}
\]
This is a concrete integer polynomial given by a finite factored expression.
It has two state ports and five local parameter ports. All auxiliary
function notation in it was explicitly defined above by polynomial
composition.

Write \(s\leadsto t\) if there exist \(a,b,c,u,v\in\N\) with
\(H(s,t;a,b,c,u,v)=0\). Because the integers have no zero divisors and
a sum of integer squares vanishes exactly when each summand does,
\(s\leadsto t\) holds exactly when one of the following occurs:
\begin{enumerate}
 \item One of the nine unconditional axioms is pushed onto \(s\).
 \item A bundle \(B(a,b,c)\) with \(c\ne a,b\) is pushed onto \(s\).
 \item The head \(b\) is replaced by \(\All(a,b)\), preserving the tail.
 \item The top two entries \(\Imp(a,b),a\), in that order, are replaced
       by \(b\), preserving the tail.
\end{enumerate}
For the bundle branch the slack witnesses exist exactly when the two
differences are nonzero: \(u=(c-a)^2-1\) and \(v=(c-b)^2-1\) are natural
precisely then. The same parameter names may be reused in different
factors, since only one factor need vanish.

\begin{theorem}[Exact accepting-path equivalence]
There is a finite \(\leadsto\)-path from 0 to 2 if and only if \(1\in D\).
\end{theorem}
\begin{proof}
Call a stack good if it is a finite list whose entries belong to \(D\).
The empty stack is good. Each push branch pushes an element of \(D\).
For a generalization or modus-ponens branch, injectivity of \(\J\)
ensures that the specified decomposition really is the top portion of
the old good stack. The corresponding rule of \(D\) proves that its
replacement head is in \(D\). Its tail remains good. Thus every state
reachable from 0 is a good finite stack. Since
\(2=\J(1,0)\), acceptance implies \(1\in D\).

Conversely, induction on a finite derivation of \(x\in D\) shows that,
for every stack \(S\), there is a path from \(S\) to \(\J(x,S)\).
For truth or an axiom use one push. For generalization, use the path
for its premise and then the generalization branch. For modus ponens,
first run the path for \(a\), then the path for \(\Imp(a,b)\), with
the resulting stack as its tail. This gives
\(\J(\Imp(a,b),\J(a,S))\), to which the MP branch applies.
Sharing a premise in a proof graph causes no problem: use a finite
proof tree, repeating its subderivation where needed. Taking \(x=1\)
and \(S=0\) proves the claim.
\end{proof}

\section{The exact ZFC endpoint}

The endpoint \path{KEndpoint.c222_iff_zf_refutation} proves
\[
 D(1)\ \longleftrightarrow\ \texttt{BProv ZFax\_s [] Form.fBot}.
\]
The endpoint \path{ConstructibleChoice.zf_refutation_iff_zfc_refutation} proves
\[
 \texttt{BProv ZFax\_s [] Form.fBot}\ \longleftrightarrow\
 \texttt{BProv ZFCax\_s [] Form.fBot}.
\]
In particular,
\[
 \text{there is a finite path }0\leadsto\cdots\leadsto2
 \quad\longleftrightarrow\quad
 \text{ZFC has a finite derivation of falsity}.
\]
Negating both sides gives the corresponding consistency assertion.
The endpoint has no consistency, transitive-model, or \(\omega\)-model
premise. Semantic constructions occur inside its proof, not as additional
assumptions of its conclusion.

\paragraph{Verification boundary.}
The numerical rows and endpoint were read from the repository sources;
their pinned Git blob hashes were checked. The new one-stack equivalence
above has a mathematical proof, but has not been compiled by Lean or Rocq
in this extraction. The repository's verification README reports external
proof-assistant results and expressly does not supply an internally encoded
PA/ZFC derivation. A claim of a theorem \emph{provable in ZFC} requires the
ordinary formalization of the syntactic, set-theoretic and Diophantine
arguments in ZFC, or a separate internal proof artifact; Lean acceptance
alone must not be relabeled as such an artifact.

\paragraph{Fixed-arity boundary.}
For a chosen \(n\), a path is expressed by \(s_0=0\), \(s_n=2\) and
\[H(s_i,s_{i+1};a_i,b_i,c_i,u_i,v_i)=0\] for \(i<n\).
The existential choice of an unbounded \(n\) leaves a variable number of
unknowns and equations. A correct fixed-arity encoding of these finite
histories must still be supplied before calling this a single Diophantine
equation for consistency.

\section{Source provenance}
All links below are at commit
\path{cdd582bec13f908df3711fa9e05d1ddac368a5db} of
\path{VladimirReshetnikov/ProveIt}. Relevant files are under
\path{Computability/BusyBeaver/Lean/BusyBeaver/ZFC/}:
\begin{itemize}
 \item \texttt{Code254.lean}: \(P,\Imp,\Rel,\All,\Eq,\Neg,\Ex\),
       the eight inherited row definitions.
 \item \texttt{KCode222.lean}: merged equality row \(g\), seed bundle
       \(B\), and the inductive calculus \texttt{C222}.
 \item \texttt{Enum222.lean}: the six-register abstract enumerator and
       its exact halting equivalence, which the direct stack construction
       does not need to simulate.
 \item \texttt{KEndpoint.lean}, \texttt{KUpper222.lean},
       \texttt{KLowerQHyp.lean}: the ZF refutation endpoint.
 \item \texttt{ConstructibleChoice.lean}: transport between finite ZF
       and ZFC refutations.
\end{itemize}
The archived source manifest records byte lengths, SHA-256 checksums and
Git blob identifiers. Sources were treated as inert data; no upstream
scripts or proof compilers were run.
\end{document}
